\section*{अध्याय \ref{ch:b2:nvs} --- मानदंडित सदिश समष्टियाँ}

\begin{solution}{exo:b2:nvs:1}
इकाई गोले: एक हीरा ($\norm\cdot_1$), एक चक्रिका ($\norm\cdot_2$), एक वर्ग ($\norm\cdot_\infty$) --- इसी क्रम
में एक के भीतर एक। असमिकाएँ: $\norm x_\infty \leq \norm x_2$ (एक वर्ग अधिकतम योग जितना होता है);
$\norm x_2 \leq \norm x_1$ (वर्ग करने पर: $x_1^2 + x_2^2 \leq (\abs{x_1}
+ \abs{x_2})^2$); $\norm x_1 \leq 2\norm x_\infty$ (दो पद, और हर एक $\leq \max$)। तीक्ष्णता: $(1, 0)$
पहली दो को समता बना देता है; $(1, 1)$ $\norm x_1 = 2\norm x_\infty$ देता है और यह भी दिखाता है कि
दूसरी दिशा में चरम अनुपात $\norm x_2 = \sqrt2 \norm x_\infty$ तथा $\norm x_1 = \sqrt2 \norm
x_2$ हैं।
\end{solution}

\begin{solution}{exo:b2:nvs:2}
\omterm{def:b2:nvs:norm}{मानदंड} अभिगृहीत: समघातता और त्रिभुज असमिका पद-दर-पद विरासत में मिलती
हैं; पृथक्करण: $N(f) = 0$ $f' = 0$ को बाध्य करता है (अतः $f$ अचर) और $f(0) = 0$: $f = 0$।
इसलिए यह \omterm{def:b2:nvs:norm}{मानदंड} है।

तुलना: $\norm f_\infty \leq N(f)$, क्योंकि $\abs{f(x)} \leq
\abs{f(0)} + \abs{\int_0^x f'} \leq \abs{f(0)} +
\norm{f'}_\infty$। विलोम विफल है: $f_n(x) = \frac1n
\sin(nx)$ लीजिए: तब $\norm{f_n}_\infty \leq \frac1n \to 0$ जबकि $N(f_n) = 0 + \norm{\cos(nx)}_\infty = 1$।
कोई अचर $C$ $N \leq C\norm\cdot_\infty$ नहीं देता।
\end{solution}

\begin{solution}{exo:b2:nvs:3}
$\abs{u(f)} \leq \norm f_\infty \int_0^1 \eu^t\,\dd t = (\eu -
1)\norm f_\infty$, जिसमें $f \equiv 1$ के लिए समता: $\vertiii u =
\eu - 1$।

स्थानांतरण: $\norm{S(x)}_\infty = \max(\abs{x_2}, \dots, \abs{x_n}) \leq
\norm x_\infty$, और $x = e_2$ पर समता: $\vertiii S = 1$ ($n
\geq 2$ के लिए)।
\end{solution}

\begin{solution}{exo:b2:nvs:4}
ऊपरी परिबंध: $\norm x_\infty \leq 1$ के लिए
\[
\abs{(Ax)_i} = \Bigl|\sum_j a_{ij}x_j\Bigr| \leq \sum_j
\abs{a_{ij}},
\]
अतः $\norm{Ax}_\infty \leq \max_i \sum_j \abs{a_{ij}}$। प्राप्ति: महत्तम को साकार करने वाला $i_0$ लीजिए और $x_j = \operatorname{sign}(a_{i_0
j})$ (जिसकी
प्रविष्टियों का मापांक $1$ हो): तब $(Ax)_{i_0} = \sum_j \abs{a_{i_0
j}}$। अतः सूत्र सिद्ध। दिए गए आव्यूह
के लिए: पंक्ति-योग $3$ और $4$: $\vertiii A_\infty = 4$।
\end{solution}

\begin{solution}{exo:b2:nvs:5}
\emph{गुणोत्तर श्रेणी:} $\vertiii B < 1$ के लिए श्रेणी $\sum
B^k$ बानाख $\mathcal{M}_n(K)$ में निरपेक्ष
रूप से अभिसरित होती है (\cref{thm:b2:nvs:absoluteconvergence}, $\vertiii{B^k} \leq
\vertiii B^k$), और गुणन की संततता (\cref{prop:b2:nvs:bilinear}) से
\[
(I - B)\sum_{k=0}^{K} B^k = I - B^{K+1} \longrightarrow I :
\]
अतः $(I - B)\sum_{k\geq0} B^k = I$, इसलिए $I - B$ प्रतिलोमनीय है और उसका प्रतिलोम वही योग है (तथा
$\vertiii{(I-B)^{-1}} \leq \frac{1}{1 - \vertiii B}$)।

\emph{विवृतता:} प्रतिलोमनीय $A$ और $\vertiii H <
\frac{1}{\vertiii{A^{-1}}}$ के लिए: $\vertiii{A^{-1}H} \leq \vertiii{A^{-1}}\vertiii H < 1$ सहित $A + H = A(I + A^{-1}H)$: अर्थात्
प्रतिलोमनीय। अतः $A$ के चारों ओर का कोई गोला $GL_n$ में ही रहता है।

\emph{प्रतिलोमन की संततता:} $B = -A^{-1}H$ के साथ
\[
(A + H)^{-1} - A^{-1} = \bigl((I - B)^{-1} - I\bigr)A^{-1}
= \Bigl(\sum_{k \geq 1} B^k\Bigr) A^{-1},
\]
जिसका \omterm{def:b2:nvs:norm}{मानदंड} $H \to 0$ होने पर $\leq \frac{\vertiii B}{1 - \vertiii B}\vertiii{A^{-1}}
\to 0$ है।
\end{solution}

\begin{solution}{exo:b2:nvs:6}
\omterm{def:b2:metric:continuity}{संतत} $\Rightarrow$ संवृत केंद्रक: संवृत $\{0\}$ का पूर्वप्रतिबिंब (\cref{thm:b2:metric:globalcontinuity})।

विलोमतः मान लीजिए $\ker\varphi$ संवृत है और $\varphi \neq 0$। $\varphi(a) = 1$ वाला $a$ चुनिए; चूँकि $a \notin \ker\varphi$
और केंद्रक संवृत है, कोई गोला $B(a, r)$ उससे नहीं मिलता। अब $\varphi(h) \neq 0$ के साथ $h \in E$
लीजिए: सदिश $a - \frac{h}{\varphi(h)}$ $\ker\varphi$ में है, अतः $B(a, r)$ के बाहर:
\[
\Bigl\Vert \frac{h}{\varphi(h)} \Bigr\Vert \geq r
\quad\Longrightarrow\quad
\abs{\varphi(h)} \leq \frac{\norm h}{r},
\]
और यह असमिका $\varphi(h) = 0$ होने पर भी तुच्छ रूप से सत्य है: अर्थात् \cref{thm:b2:nvs:continuouslinear} का
परिबंध (4): अतः \omterm{def:b2:metric:continuity}{संतत}।
\end{solution}

\begin{solution}{exo:b2:nvs:7}
मान लीजिए $f_n$ $\intcc{0}{\frac12 - \frac1n}$ पर $0$ है, फिर $\frac12$ पर मान $1$ तक आनत, और $\intcc{\frac12}{1}$ पर $1$।
$m
\geq n$ के लिए $f_m - f_n$ लंबाई $\frac1n$ के किसी अंतराल पर आधारित है और उसके मान $\intcc{-1}{1}$
में हैं: अतः $\norm{f_m - f_n}_1 \leq
\frac1n$: कोशी।

मान लीजिए $\norm\cdot_1$ में $f_n \to f$ है और $f$ \omterm{def:b2:metric:continuity}{संतत} है। $\intcc{0}{\frac12 - \delta}$ पर (स्थिर $\delta$): $n >
\frac1\delta$ के
लिए $\int \abs{f} =
\int\abs{f - f_n} \leq \norm{f - f_n}_1 \to 0$, अतः $\int_0^{1/2 - \delta}\abs f = 0$, और यथार्थ धनात्मकता से वहाँ $f = 0$ --- और यह प्रत्येक
$\delta$ के लिए: अतः $\intoo{0}{\frac12}$ पर $f = 0$। इसी प्रकार $\intcc{\frac12}{1}$ पर $f = 1$ ($f_n$ वहाँ सब
$1$ के बराबर हैं)। $\frac12$ पर संततता से: $0 =
1$, जो असंगत है। कोई सीमा नहीं
है: अतः समष्टि \omterm{def:b2:metric:complete}{पूर्ण} नहीं।
\end{solution}

\begin{solution}{exo:b2:nvs:8}
सुपरिभाषित और \omterm{def:b2:metric:continuity}{संतत}: $\abs{\varphi(f)} \leq \sum 2^{-n}
\norm f_\infty = \norm f_\infty$, अतः $\varphi$ ऐसा रैखिक फलनिक है कि $\vertiii\varphi \leq 1$ (हर $f$ के
लिए श्रेणी निरपेक्ष रूप से अभिसरित होती है)।

\omterm{def:b2:nvs:norm}{मानदंड} $1$: $K$ स्थिर कीजिए; बिंदु $1, \frac12, \dots, \frac1K$ परस्पर भिन्न हैं, अतः ऐसा \omterm{def:b2:metric:continuity}{संतत}
$f_K$, $\norm{f_K}_\infty
\leq 1$ है कि $n \leq K$ के लिए $f_K(\frac1n) = (-1)^n$ (खंडशः आनत अंतर्वेशन, $0$ के पास अचर)।
तब
\[
\varphi(f_K) \geq \sum_{n=1}^{K} 2^{-n} - \sum_{n > K} 2^{-n}
= 1 - 2^{-K+1} \xrightarrow[K \to \infty]{} 1 .
\]

प्राप्त नहीं: यदि $\norm f_\infty \leq 1$ और $\varphi(f) = 1$ हो, तो हर पद को अपना महत्तम योगदान देना
पड़ेगा: अर्थात् प्रत्येक $n$ के लिए $(-1)^n f(\frac1n) = 1$ (अन्यथा एक पद की यथार्थ कमी की
भरपाई नहीं हो सकती, क्योंकि सारे पद $\leq 2^{-n}$ हैं)। अतः $f(\frac1n) = (-1)^n$; परंतु $\frac1n \to 0$ और $f$
$0$ पर \omterm{def:b2:metric:continuity}{संतत} है, जो $(-1)^n$ के विरोधाभासी अभिसरण को बाध्य कर देता है। अतः
उच्चतम महत्तम नहीं है --- जो परिमित विमा में असंभव है, जहाँ संवृत इकाई
गोला \omterm{def:b2:metric:compact}{संहत} होता है।
\end{solution}

\begin{solution}{exo:b2:nvs:9}
\emph{परिबद्ध अनुक्रम:} कोई परिबद्ध अनुक्रम किसी संवृत गोले $\overline B(0, R) = R\,\overline B(0,1)$ में
पड़ता है, जो \omterm{def:b2:metric:compact}{संहत} है (समाकृतिकता $x \mapsto Rx$ के अंतर्गत \omterm{def:b2:metric:compact}{संहत} इकाई गोले का
प्रतिबिंब): वहीं उपअनुक्रम निकालिए।

\emph{रैखिक फलनिक:} यदि $\dim E < \infty$, तो $E$ से जाने वाला हर रैखिक प्रतिचित्रण
\omterm{def:b2:metric:continuity}{संतत} है (\cref{thm:b2:nvs:finitedim})। विलोमतः मान लीजिए $\dim E = \infty$ (जिसे \cref{thm:b2:nvs:riesz} के अनुसार संहतता वाली
परिकल्पना वास्तव में अपवर्जित कर देती है --- इस प्रश्न का उद्देश्य
सामान्य मानदंडित समष्टियों में दोनों गुणधर्मों के बीच का निहितार्थ है):
कोई सामान्यीकृत रैखिकतः स्वतंत्र अनुक्रम $(e_n)$ चुनिए, उसे बीजगणितीय आधार
तक पूरा कीजिए (स्वीकृत), और $\varphi(e_n) = n$ परिभाषित कीजिए तथा शेष आधार सदिशों पर
$\varphi = 0$, और फिर रैखिक रूप से विस्तार कीजिए। तब $\norm{e_n} = 1$ के साथ $\abs{\varphi(e_n)} = n$: अर्थात्
इकाई गोले पर अपरिबद्ध, इसलिए असंतत। अतः ``सारे फलनिक \omterm{def:b2:metric:continuity}{संतत}'' परिमित विमा
को बाध्य कर देता है।
\end{solution}

\begin{solution}{exo:b2:nvs:10}
अचर फलन $1$ के साथ कोशी--श्वार्ज़: $\norm f_1 =
\int_0^1 \abs f\cdot 1 \leq \bigl(\int_0^1
f^2\bigr)^{1/2}\bigl(\int_0^1 1\bigr)^{1/2} = \norm f_2$। और $\norm f_2^2 = \int f^2 \leq \norm f_\infty^2$। $f_n(x) =
x^n$ के लिए:
\[
\norm{f_n}_1 = \frac1{n+1}, \qquad
\norm{f_n}_2 = \frac1{\sqrt{2n+1}}, \qquad
\norm{f_n}_\infty = 1 .
\]
तब $\norm{f_n}_2/\norm{f_n}_1 = \frac{n+1}{\sqrt{2n+1}} \to
\infty$ और $\norm{f_n}_\infty/\norm{f_n}_2 = \sqrt{2n+1} \to
\infty$: अतः कोई उलटी असमिका नहीं, और कोई भी युग्म \omterm{def:b2:nvs:equivalent}{तुल्य} नहीं।
\end{solution}

\begin{solution}{exo:b2:nvs:11}
\emph{दूरी के लिए निचला परिबंध:} $h \in \ker\varphi$ के लिए $\abs{\varphi(x)} = \abs{\varphi(x - h)} \leq
\vertiii\varphi\,\norm{x - h}$; $h$ पर निम्नतम लीजिए:
$d(x, \ker\varphi) \geq \abs{\varphi(x)}/\vertiii\varphi$।

\emph{ऊपरी परिबंध:} हम $\varphi(x) \neq 0$ मान सकते हैं। $\varepsilon > 0$ दिया हो तो $\abs{\varphi(u)} \geq
\vertiii\varphi - \varepsilon > 0$ वाला कोई
इकाई $u$ चुनिए और $h = x -
\frac{\varphi(x)}{\varphi(u)}\,u$ रखिए: तब $\varphi(h) = 0$ और
\[
\norm{x - h} = \frac{\abs{\varphi(x)}}{\abs{\varphi(u)}}
\leq \frac{\abs{\varphi(x)}}{\vertiii\varphi - \varepsilon}.
\]
$\varepsilon \to 0$ लीजिए: $d(x, \ker\varphi) \leq
\abs{\varphi(x)}/\vertiii\varphi$; अर्थात् समता।

\emph{अप्राप्ति:} \cref{exo:b2:nvs:8} से $\varphi$ लीजिए ($\vertiii\varphi = 1$, जो प्राप्त नहीं होता) और $\varphi(x) \neq 0$
वाला कोई भी $x$। यदि कोई $h \in \ker\varphi$ $\norm{x - h} = \abs{\varphi(x)}$ को साकार करता, तो इकाई सदिश $v = (x -
h)/\norm{x - h}$ $\abs{\varphi(v)} =
\abs{\varphi(x)}/\norm{x - h} = 1 = \vertiii\varphi$
पूरा करता: अर्थात् \omterm{thm:b2:nvs:continuouslinear}{संकारक मानदंड} प्राप्त हो जाता --- विरोधाभास।
\end{solution}

\begin{solution}{exo:b2:nvs:12}
$N_1 \leq N_2$ प्रांत में उच्चतम की एकदिष्टता ही है, अतः तत्समक $(E, N_2) \to (E, N_1)$ $1$-लिप्शिट्स
है। $P_n(x) =
(x/2)^n$ के लिए: $N_2(P_n) = 1$ ($x = 2$ पर प्राप्त) जबकि $N_1(P_n) =
2^{-n}$: कोई परिबंध $N_2 \leq CN_1$ सभी $n$
के लिए $1 \leq C2^{-n}$ दे देता: जो असंभव है। अतः \omterm{def:b2:nvs:norm}{मानदंड} \omterm{def:b2:nvs:equivalent}{तुल्य} नहीं हैं और प्रतिलोम
तत्समक असंतत रैखिक एकैकी आच्छादक प्रतिचित्रण है।

\emph{अपूर्णता:} $S_n = \sum_{k=0}^{n}\frac{X^k}{k!}$ लीजिए। $m > n$ के लिए $N_1(S_m - S_n) \leq \sum_{k>n}\frac1{k!} \to 0$: अर्थात् $N_1$ के लिए कोशी।
यदि $(E, N_1)$ में $S_n \to P$ हो, तो बिंदुवार $\intcc{0}{1}$ पर $P(x) = \lim S_n(x) = \eu^x$; परंतु घात $d$ का कोई बहुपद
किसी अंतराल पर $\eu^x$ के बराबर नहीं हो सकता ($d + 1$ बार अवकलन कीजिए: बायाँ
पक्ष मर जाता है, $\eu^x$ नहीं)। $E$ में कोई सीमा नहीं: अतः \omterm{def:b2:metric:complete}{पूर्ण} नहीं।

परिमित विमा में तीनों परिघटनाएँ असंभव हैं: सारे \omterm{def:b2:nvs:norm}{मानदंड} \omterm{def:b2:nvs:equivalent}{तुल्य} होते हैं,
हर मानदंडित समष्टि \omterm{def:b2:metric:complete}{पूर्ण} होती है, और रैखिक एकैकी आच्छादक प्रतिचित्रण का
प्रतिलोम परिमित-विमीय समष्टि से जाने वाला रैखिक प्रतिचित्रण है, अतः \omterm{def:b2:metric:continuity}{संतत}
(\cref{thm:b2:nvs:finitedim})।
\end{solution}

\begin{solution}{pb:b2:nvs:1}
\textbf{1.} विचार करने योग्य उम्मीदवार $K = \{f \in F
: \norm{x - f} \leq \norm x\}$ बनाते हैं: यह अरिक्त ($0 \in K$),
संवृत (\omterm{def:b2:metric:continuity}{संतत} $f \mapsto
\norm{x - f}$ के अंतर्गत किसी संवृत अंतराल का पूर्वप्रतिबिंब, संवृत
$F$ से प्रतिच्छेदित, \cref{cor:b2:nvs:closedsubspace}) और परिबद्ध ($\norm f \leq
\norm{f - x} + \norm x \leq 2\norm x$) है। परिमित-विमीय $F$ में
संवृत और परिबद्ध का अर्थ \omterm{def:b2:metric:compact}{संहत} है (\cref{thm:b2:nvs:finitedim}); और \omterm{def:b2:metric:continuity}{संतत} फलन $f
\mapsto \norm{x - f}$ $K$ पर अपना
निम्नतम प्राप्त कर लेता है, जो पूरे $F$ पर लिए गए निम्नतम के बराबर है
(कोई भी $f \notin K$ $\norm{x -
f} > \norm x \geq \inf$ देता है)।

\textbf{2.} $(\R^n,
\norm\cdot_2)$ में समांतर चतुर्भुज सर्वसमिका: $\norm{u + v}^2 + \norm{u - v}^2 = 2\norm u^2 +
2\norm v^2$ (निर्देशांकों के
योगों के वर्ग खोलिए)। इकाई $u \neq v$ के लिए:
\[
\Bigl\Vert\frac{u+v}2\Bigr\Vert^2 = 1 - \frac{\norm{u -
v}^2}{4} < 1 .
\]
यथार्थतः उत्तल नहीं: $\norm\cdot_\infty$ के लिए $u = (1, 1, 0,
\dots)$, $v = (1, -1, 0, \dots)$ लीजिए: ये इकाई सदिश हैं जिनका
मध्यबिंदु $(1, 0, \dots)$ है और \omterm{def:b2:nvs:norm}{मानदंड} $1$; और $\norm\cdot_1$ के लिए $u = (1,
0, \dots)$, $v = (0, 1, 0, \dots)$ लीजिए:
मध्यबिंदु $(\frac12,
\frac12, 0, \dots)$, \omterm{def:b2:nvs:norm}{मानदंड} $1$।

\textbf{3.} मान लीजिए $d = d(x, F)$। यदि $d = 0$: तो $x \in \overline F =
F$ और एकमात्र न्यूनतमकारी $x$
है। और यदि $d > 0$ तथा $f_1 \neq f_2$ दोनों न्यूनतम करें: तो $u = \frac{x - f_1}{d}$ और $v = \frac{x -
f_2}{d}$ भिन्न इकाई
सदिश हैं, अतः $\frac{f_1 + f_2}2 \in F$ सहित
\[
\Bigl\Vert x - \frac{f_1 + f_2}2\Bigr\Vert
= d\,\Bigl\Vert\frac{u + v}2\Bigr\Vert < d ,
\]
जो $d$ की परिभाषा के विरुद्ध है। इसलिए न्यूनतमकारी अद्वितीय है।

\textbf{4.} $\norm{(0,1) - t(1,0)}_\infty = \max(\abs t, 1)
\geq 1$, जिसमें समता तभी जब $\abs t \leq 1$: अतः न्यूनतमकारी खंड $\{t(1, 0) : t \in \intcc{-1}{1}\}$ बनाते
हैं, जो सब दूरी $1$ पर हैं --- और अद्वितीयता ठीक इसलिए विफल होती है कि
वर्गाकार गोले की भुजाएँ सपाट हैं (प्रश्न 2)।

\textbf{5.} $M = \max f$, $m = \min f$ लीजिए (जो संहतता से प्राप्त होते हैं)। किसी अचर
$c$ के लिए: $\sup\abs{f - c} \geq
\max(M - c,\, c - m) \geq \frac{M - m}2$, जहाँ अंतिम असमिका इसलिए है कि दोनों राशियों का औसत $\frac{M-m}2$
है; और दोनों में समता $M - c = c - m$ को बाध्य कर देती है, अर्थात् $c = c^* = \frac{M + m}2$। विलोमतः
$\sup\abs{f - c^*} = \max(M - c^*, c^* - m) =
\frac{M - m}2$। अतः सर्वोत्तम अचर अद्वितीय है। $\intcc01$ पर $f(x) = x^2$ के लिए: $c^* = \frac12$, दूरी $\frac12$।

\textbf{6.} $\cos(n{+}1)\theta + \cos(n{-}1)\theta =
2\cos\theta\cos n\theta$ से: $T_0 = 1$, $T_1 = X$, $T_{n+1} = 2XT_n - T_{n-1}$ से परिभाषित बहुपद आगमन से $T_n(\cos\theta)
= \cos n\theta$ पूरा
करते हैं। अद्वितीयता: $\intcc{-1}{1}$ (अपरिमित बिंदुओं) पर सहमत दो बहुपद बराबर होते
हैं। फिर आगमन: $n \geq 1$ के लिए अग्र गुणांक $2^{n-1}$ के साथ $\deg T_n = n$ ($T_1$: गुणांक $1 = 2^0$;
और पुनरावृत्ति उसे दुगुना कर देती है)।

\textbf{7.} हर $x \in \intcc{-1}{1}$ $\cos\theta$ है, और $\abs{\cos n\theta} \leq 1$। $\eta_k = \cos\frac{k\pi}n$ पर: $T_n(\eta_k) = \cos k\pi = (-1)^k$, और $1 = \eta_0 > \eta_1 >
\dots > \eta_n = -1$। मूल: $\cos n\theta = 0$ यदि
और केवल यदि $\theta =
\frac{(2k-1)\pi}{2n}$: अर्थात् $n$ भिन्न बिंदु $\cos\frac{(2k-1)\pi}{2n}$, और चूँकि $\frac{(k-1)\pi}n <
\frac{(2k-1)\pi}{2n} < \frac{k\pi}n$, हर मूल दो
क्रमागत चरम बिंदुओं के ठीक बीच में पड़ता है।

\textbf{8.} $T_2 = 2X^2 - 1$, $T_3 = 4X^3 - 3X$, $T_4 = 8X^4 -
8X^2 + 1$। $T_3$ के लिए: $T_3(1) = 1$, $T_3(\tfrac12) = \tfrac12 -
\tfrac32 = -1$, $T_3(-\tfrac12) = 1$, $T_3(-1) = -1$: अर्थात्
पूर्ण एकांतरण।

\textbf{9.} $x \geq 1$ के लिए $u_\pm = x \pm \sqrt{x^2 - 1}$ लीजिए: ये $z^2 - 2xz + 1$ के मूल हैं, जहाँ $u_+u_- = 1$। अनुक्रम
$s_n = \frac{u_+^n + u_-^n}2$ $s_{n+1} = 2x\,s_n -
s_{n-1}$ पूरा करता है (द्विघात से न्यूटन-प्रकार की पुनरावृत्ति), $s_0 = 1$,
$s_1 = x$: अर्थात् $n \mapsto
T_n(x)$ जैसी ही पुनरावृत्ति और आरंभिक मान, अतः सभी $n$ के लिए
$s_n = T_n(x)$। चूँकि $x > 1$ के लिए $u_+ > 1$ सहित $0 < u_- \leq 1
\leq u_+$: $T_n(x) \geq
\frac{u_+^n}2 \to \infty$ और $T_n(x) \sim \frac12\bigl(x +
\sqrt{x^2-1}\bigr)^n$। ($x \leq -1$ के लिए
सम-विषमता $T_n(-x)
= (-1)^nT_n(x)$ का उपयोग कीजिए, जो पुनरावृत्ति से स्पष्ट है।)

\textbf{10.} प्रत्येक $\theta$ के लिए: $T_m\bigl(T_n(\cos\theta)\bigr)
= T_m(\cos n\theta) = \cos mn\theta =
T_{mn}(\cos\theta)$। बहुपद $T_m \circ T_n$ और $T_{mn}$ $\intcc{-1}{1}$ पर सहमत
हैं, अतः बराबर हैं।

\textbf{11.} $D = Q_n - P$ की घात $\leq n - 1$ है (एकिक अग्र पद कट जाते हैं)। चरम बिंदुओं
पर: परिकल्पना से $(-1)^kD(\eta_k) =
2^{1-n} - (-1)^kP(\eta_k) \geq 2^{1-n} - \abs{P(\eta_k)} > 0$। अतः $D$ $n + 1$ ह्रासमान बिंदुओं $\eta_0 > \dots > \eta_n$ पर बारी-बारी
चिह्न वाले अशून्य मान लेता है: मध्यमान प्रमेय से उसके कम से कम $n$
भिन्न मूल हैं, प्रत्येक \omterm{def:b2:metric:topology}{विवृत} अंतराल $\intoo{\eta_{k}}{\eta_{k-1}}$ में एक। घात $\leq n - 1$ का अशून्य बहुपद
$n$ मूल नहीं रख सकता; और $D = 0$ यथार्थ चिह्नों के विरुद्ध है। अतः
विरोधाभास: घात $n$ के प्रत्येक एकिक $P$ के लिए $\sup\abs P \geq 2^{1-n}$।

\textbf{12.} अब $\abs{P} \leq 2^{1-n}$ के कारण $(-1)^kD(\eta_k) = 2^{1-n} - (-1)^kP(\eta_k)
\geq 0$। प्रत्येक $\intcc{\eta_k}{\eta_{k-1}}$ ($k = 1, \dots, n$) पर $D$ के
अंत्यबिंदु-मानों के दुर्बल चिह्न विपरीत हैं: मध्यमान प्रमेय संवृत
अंतराल में कोई शून्य $z_k$ दे देती है। यदि $z_k$ परस्पर भिन्न चुने जा सकें,
तो घात $\leq n-1$ वाले $D \neq 0$ के $n$ मूल होंगे: विरोधाभास। दो क्रमागत अंतराल
केवल शून्य $z_k = z_{k+1} = \eta_k$ साझा कर सकते हैं, जहाँ $0 < k < n$ (अभ्यंतरीय)। वहाँ $D(\eta_k) = 0$ का
अर्थ है $P(\eta_k) =
(-1)^k2^{1-n}$, जो $\intcc{-1}{1}$ पर $P$ का चरम मान है और किसी \emph{अभ्यंतरीय}
बिंदु पर प्राप्त: अतः $P'(\eta_k) = 0$; और $\eta_k$ $T_n$ का भी अभ्यंतरीय चरम है: अतः
$Q_n'(\eta_k) =
0$। इसलिए $D'(\eta_k) = 0$: $\eta_k$ बहुलता $\geq 2$ का मूल है, जो साझे अंतराल की भरपाई
कर देता है। सभी स्थितियों में $D$ के बहुलता सहित कम से कम $n$ मूल हैं
जबकि घात $\leq n -
1$ है, अतः $D = 0$: $P = Q_n$। चेबिशेव की चरम प्रमेय सिद्ध हुई:
अद्वितीय एकिक न्यूनतमकारी $2^{1-n}T_n$ है, जिसका उच्चतम \omterm{def:b2:nvs:norm}{मानदंड} $2^{1-n}$ है।

\textbf{13.} घात $n$ के एकिक बहुपद ठीक $R \in \R_{n-1}[X]$ वाले $X^n - R$ हैं, अतः
\[
d_\infty\bigl(X^n, \R_{n-1}[X]\bigr) = \min_{P \text{
एकिक}}\ \sup_{\intcc{-1}{1}}\abs P = 2^{1-n},
\]
जो अद्वितीय रूप से $R^* = X^n - Q_n$ पर होता है। प्रतिस्थापन $x = \frac{1+t}2$: यदि $P$ $\intcc01$ पर
घात $n$ का एकिक है, तो $t \mapsto
2^nP\bigl(\frac{1+t}2\bigr)$ $\intcc{-1}{1}$ पर एकिक है और उसका उच्चतम $2^n\sup_{\intcc01}\abs P$ के
बराबर है: अतः $\sup_{\intcc01}\abs P \geq 2^{-n}\cdot2^{1-n} = 2^{1-2n}$, जिसमें समता ठीक $P^*(x) = 2^{-n}Q_n(2x - 1)$ के लिए: इसलिए $\intcc{0}{1}$ पर दूरी
$2^{1-2n}$ है।

\textbf{14.} $\omega$ घात $n$ का एकिक है, अतः $\sup_{\intcc{-1}{1}}\abs\omega \geq 2^{1-n}$, जिसमें समता तभी जब $\omega = Q_n = 2^{1-n}T_n$,
अर्थात् तभी जब नोड $T_n$ के $n$ मूल हों: $x_k = \cos\frac{(2k-1)\pi}{2n}$ --- अर्थात् \emph{चेबिशेव
नोड}। न्यूनतम मान: $2^{1-n}$। समदूरस्थ नोड यथार्थतः बदतर हैं; अंतर्वेशन
त्रुटि-गुणक नोडों को अंत्यबिंदुओं के पास गुच्छित करने से न्यूनतम होता
है।

\textbf{15.} हाथ से $n = 2$: $x^2$ $\intcc01$ पर घूमता है, अतः $\sup_{\intcc{-1}{1}}\abs{x^2 - c} = \max(\abs c, \abs{1 - c})
\geq \frac12$, जो $c = \frac12$ पर
न्यूनतम होता है: अर्थात् न्यूनतम एकिक द्विघात $x^2 - \frac12 = \frac12(2x^2 - 1) = Q_2$, मान $\frac12 = 2^{1-2}$। $\intcc{0}{1}$ पर $n = 10$
के लिए: $2^{1-20} =
2^{-19} \approx 1.9\cdot10^{-6}$। घात $9$ का कोई बहुपद पूरे $\intcc01$ पर $x^{10}$ के दो लाखवें भाग
के भीतर बना रहता है: इस पैमाने पर दोनों आलेख अभेद्य हैं --- $0$ के पास
$x^{10}$ का चपटापन कम घातों से ही सारा काम करा लेता है।

\textbf{16.} $\norm{g_k}_\infty = g_k(1) = 1$। $j > k$ के लिए $a = 2^k$, $b = 2^j \geq 2a$ रखिए और $x_0 = 2^{-1/a}$ पर मूल्यांकन कीजिए
(अतः $x_0^a = \frac12$):
\[
g_k(x_0) - g_j(x_0) = \frac12 - \Bigl(\frac12\Bigr)^{b/a}
\geq \frac12 - \frac14 = \frac14 .
\]
अतः सभी $j \neq k$ के लिए $\norm{g_k - g_j}_\infty \geq \frac14$: कोई उपअनुक्रम कोशी नहीं, इसलिए कोई अभिसरित भी
नहीं। $\bigl(C(\intcc01), \norm\cdot_\infty\bigr)$ का संवृत इकाई गोला \omterm{def:b2:metric:compact}{संहत} नहीं है।

\textbf{17.} $F$ संवृत (\cref{cor:b2:nvs:closedsubspace}) और उचित है: $y \notin F$ चुनिए, अतः $\delta = d(y, F) > 0$। प्रश्न 1
से दूरी किसी $f^* \in F$ पर प्राप्त होती है। $x
= \frac{y - f^*}{\delta}$ रखिए, जो इकाई सदिश है ($\norm{y - f^*} =
\delta$)।
प्रत्येक $g \in F$ के लिए:
\[
\norm{x - g} = \frac{\norm{y - (f^* + \delta g)}}{\delta}
\geq \frac{\delta}{\delta} = 1 ,
\]
क्योंकि $f^* + \delta g \in F$। अतः $d(x, F) \geq 1$; और $d(x,
F) \leq \norm{x - 0} = 1$: अर्थात् ठीक $1$।

\textbf{18.} अनंत-विमीय $E$ में आगमन से इकाई सदिश बनाइए: $x_1$ स्वेच्छ;
और $x_1, \dots, x_k$ दिए हों तो उपसमष्टि $F_k = \operatorname{Vect}(x_1, \dots, x_k)$ परिमित-विमीय है, अतः उचित, और प्रश्न 17
$d(x_{k+1}, F_k) = 1$ वाला इकाई $x_{k+1}$ दे देता है: विशेष रूप से $i \leq k$ के लिए $\norm{x_{k+1} - x_i} \geq 1$। इस अनुक्रम
की परस्पर दूरियाँ $\geq 1$ हैं: अतः इकाई गोले में बिना अभिसारी उपअनुक्रम का
अनुक्रम है, इसलिए वह \omterm{def:b2:metric:compact}{संहत} नहीं --- यही रीस की प्रमेय है, तीखे अचर $1$
के साथ।

\textbf{19.} मान लीजिए $h_n$ $\bigl[\frac1{n+1}, \frac1n\bigr]$ के दोनों आधों पर आनत है, $0$ से मध्यबिंदु
पर $1$ तक और फिर $0$ तक, और अन्यत्र शून्य: यह \omterm{def:b2:metric:continuity}{संतत} है और $\norm{h_n}_\infty = 1$। $n \neq m$ के
लिए आधार अधिक से अधिक किसी साझे अंत्यबिंदु पर मिलते हैं, जहाँ दोनों
लुप्त हैं; और $h_n$ के शिखर पर $h_m = 0$: अतः ठीक $\norm{h_n - h_m}_\infty = 1$। स्थिर $x > 0$ के लिए: $\frac1n < x$
होते ही $h_n(x) = 0$, और सदा $h_n(0) = 0$: अतः बिंदुवार $h_n \to 0$; परंतु $\norm{h_n - 0}_\infty = 1$: एकसमान नहीं।
यही इकाई गोले के भीतर परस्पर दूरी $1$ वाला स्पष्ट नक्षत्र है।

\textbf{20.} त्रिज्या $\frac13$ के गोले का व्यास $\leq
\frac23 < 1$ है, अतः उसमें $h_n$ में से
अधिकतम एक ही आ सकता है (दो की दूरी $1$ है)। ऐसे परिमित गोलों में
अपरिमित $h_n$ में से परिमित ही आते हैं: अतः वे इकाई गोले को आच्छादित नहीं
कर सकते। पूर्ण परिबद्धता --- जो \cref{thm:b2:metric:borellebesgue} की उपपत्ति के अनुसार \omterm{def:b2:metric:compact}{संहत} दूरिक
समष्टियों को प्राप्त है --- यहाँ जितनी बुरी तरह विफल हो सकती थी उतनी
होती है।

\textbf{21.} यदि $x \neq 0$ के साथ $Ax = \lambda x$: तो $\abs\lambda\,\norm x = \norm{Ax} \leq \vertiii A\,\norm x$, अतः $\abs\lambda \leq \vertiii A$। $\begin{psmallmatrix}1 & -2\\ 3 & 1\end{psmallmatrix}$ के लिए:
$\vertiii A_\infty = \max(1 + 2,\ 3 + 1) = 4$
(\cref{exo:b2:nvs:4}), अतः हर \omterm{def:b2:reduction:eigen}{अभिलक्षणिक मान} का मापांक $\leq
4$ है; वास्तव में $\chi_A = X^2 - 2X + 7$ मापांक
$\sqrt7 \approx 2.65$ वाले $\lambda = 1 \pm
\iu\sqrt6$ देता है: अर्थात् परिबंध मान्य है, तीखा नहीं।

\textbf{22.} $N_P$ एक \omterm{def:b2:nvs:norm}{मानदंड} है: $N_P(x) = 0$ $P^{-1}x = 0$ को बाध्य करता है, अतः $x = 0$; और
समघातता तथा त्रिभुज असमिका रैखिक $P^{-1}$ के द्वारा $\norm\cdot_\infty$ से विरासत में मिलती
हैं। \omterm{thm:b2:nvs:continuouslinear}{संकारक मानदंड}: $y = P^{-1}x$ और $D = \mathrm{diag}(\lambda_i)$ के साथ
\[
N_P(Ax) = \norm{P^{-1}AP\,y}_\infty = \norm{Dy}_\infty,
\]
अतः $\vertiii A_{N_P}$ $D$ का $\norm\cdot_\infty$-संकारक \omterm{def:b2:nvs:norm}{मानदंड} है, जो उसका सबसे बड़ा निरपेक्ष
पंक्ति-योग है (\cref{exo:b2:nvs:4}): $\max_i\abs{\lambda_i}$।

\textbf{23.} यदि सभी $\abs{\lambda_i} < 1$: तो $\rho = \max\abs{\lambda_i} < 1$ के साथ $N_P(A^kx) \leq
\rho^kN_P(x)$, अतः प्रत्येक $x$ के लिए
$A^kx
\to 0$, और $\mathcal{M}_n$ के किसी भी \omterm{def:b2:nvs:norm}{मानदंड} में $A^k \to 0$ (परिमित विमा में सब \omterm{def:b2:nvs:equivalent}{तुल्य} हैं,
\cref{thm:b2:nvs:finitedim}; और हर $j$ के लिए $A^ke_j$ का अभिसरण प्रविष्टि-दर-प्रविष्टि अभिसरण है)।
और यदि \omterm{def:b2:reduction:eigen}{अभिलक्षणिक सदिश} $x$ के साथ कोई $\abs{\lambda} \geq 1$: तो $\norm{A^kx} = \abs\lambda^k\norm x
\not\to 0$। $A = \frac14\begin{psmallmatrix}1 & 2\\ 2 &
1\end{psmallmatrix}$ के लिए:
\omterm{def:b2:reduction:eigen}{अभिलक्षणिक मान} $\frac14(1 \pm 2) =
\frac34, -\frac14$, दोनों का मापांक $< 1$: अतः $A^k \to 0$।

\textbf{24.} दोनों परिमित-विमीय $\R_n[X]$ पर \omterm{def:b2:nvs:norm}{मानदंड} हैं, अतः हर $n$ के लिए
\omterm{def:b2:nvs:equivalent}{तुल्य}। $\intcc01$ पर प्रश्न 13 का न्यूनतम एकिक $P^*_n$ लीजिए: $X^n$ का उसका गुणांक
$1$ है, अतः $N_c(P^*_n) \geq 1$, जबकि $\norm{P^*_n}_{\intcc01} = 2^{1-2n}$। इसलिए
\[
C_n \geq \frac{N_c(P^*_n)}{\norm{P^*_n}_{\intcc01}}
\geq 2^{2n-1} .
\]
तुल्यता के अचर विमा के साथ फट जाते हैं: सम्मिलन $\R[X]$ पर कोई एक अचर काम
नहीं आता, और यही \cref{exo:b2:nvs:12} में देखी गई अतुल्यता है --- ``सारे \omterm{def:b2:nvs:norm}{मानदंड} \omterm{def:b2:nvs:equivalent}{तुल्य}
हैं'' एक बार में एक ही विमा के बारे में प्रमेय है, और अनंत विमा वह जगह
है जहाँ वह मर जाती है।

\textbf{25.} परिमित विमा में संवृत परिबद्ध समुच्चयों की संहतता ने
सर्वोत्तम सन्निकटनों का अस्तित्व दिया (प्रश्न 1) और $n + 1$ चरम बिंदुओं पर
शून्य-गणना को चलाया (प्रश्न 11--12, प्राप्त उच्चतमों के द्वारा)।
यथार्थ उत्तलता सर्वोत्तम सन्निकटन की अद्वितीयता पर शासन करती है ---
गोल गोले एक ही न्यूनतमकारी देते हैं, सपाट भुजाओं वाले गोले उनके पूरे
खंड (प्रश्न 2--4)। परस्पर दूरी $1$ वाले इकाई सदिशों का नक्षत्र (प्रश्न
17--19) इकाई गोले की संहतता और उसके साथ पूर्ण परिबद्धता भी नष्ट कर देता
है (प्रश्न 20)। प्रश्न 24 उस ढहने को परिमाणित करता है: $\R_n[X]$ पर दो
मानदंडों को जोड़ने वाले अचर $4^n$ की तरह बढ़ते हैं, अतः $\R[X]$ तक जाते-जाते
कोई एकसमान तुलना नहीं बचती। शिखर चेबिशेव की चरम प्रमेय है (प्रश्न
11--12): अद्वितीय एकिक न्यूनतमकारी $2^{1-n}T_n$। और सर्वोत्तम सन्निकटन को अपना
आधुनिक घर हिल्बर्ट समष्टियों में मिलता है, जहाँ प्रक्षेप प्रमेय संहतता
के स्थान पर पूर्णता और समांतर चतुर्भुज सर्वसमिका ले आती है --- जिसे
तृतीय वर्ष के खंड में ईमानदारी से सिद्ध किया गया है।
\end{solution}
